% !PRINTED 453
\weberpage{453}
%<<CONT
\weq{15}{p\,\omega, \quad \frac{\omega+c}{p}}
no two equivalent ones, then to each value of $u$ there belong $[p-(D,p)]$ values $v$, and if we denote the degrees of (2) and (3) by $h$ and $h'$, then the following relation results (in agreement with the formulae of \S\,100):

\weq{16}{h' = [p-(D,p)]\,h.}

According to \S\,122, $h$ and $h'$ are the class numbers of the discriminants $D, D'$.

It thus remains to ask whether, among the quantities (15), after the exclusion of those which fall away by reason of (13), equivalent ones can still occur. If in the first place $p\omega$ is equivalent to $\frac{\omega+c}{p}$, then:

\weq{17}{\frac{\omega+c}{p} = \frac{\gamma+\delta p\omega}{\alpha+\beta p\omega}, \quad \alpha\delta-\beta\gamma=1,}

or:

\weq{18}{\beta p\,\omega^2 + (\alpha+\beta c p - \delta p^2)\omega + c\alpha - \gamma p = 0.}

This equation must arise from (9) by multiplication with an integer $g$, and it therefore follows that:
\[
\beta p = gA, \quad \alpha+\beta c p-\delta p^2 = gB, \quad c\alpha - gp = gC. % [sic] German print has g p; elimination from (18) requires \gamma p
\]
Since $A$ was assumed to be not divisible by $p$, it follows from this that $g$, and consequently $\alpha$, is divisible by $p$. If we set $g=px$, $\alpha=p\alpha'$, then it follows that:

\weq{19}{\beta=Ax, \quad \alpha'+\beta c-\delta p=Bx, \quad c\alpha'-\gamma=Cx.}

If we further set

\weq{20}{\alpha'-\beta c+\delta p=y,}

then we obtain:

\weq{21}{\begin{aligned} 2\alpha' &= Bx+y,\\ 2(\beta c-\delta p) &= Bx-y, \end{aligned}}

and from

\weq{22}{p\alpha'\delta-\beta\gamma=1}

there then follows

\weq{23}{y^2-Dx^2=4.}

The value $x=0$ is to be excluded, since otherwise $\beta=0$ would have to hold, which contradicts equation (22), and we may, without loss of generality, assume $x$ to be positive. Since moreover $D$ is, with respect to the modulus 4, either $\equiv 0$ or $\equiv 1$, then
%CONT>>
